\documentclass[11pt,a4paper]{article}
\usepackage[utf8]{inputenc}
\usepackage[T1]{fontenc}
\usepackage{amsmath,amssymb,amsthm}
\usepackage{listings}
\usepackage{geometry}
\usepackage{hyperref}
\geometry{margin=2.5cm}

\newtheorem{theorem}{Theorem}
\newtheorem{lemma}[theorem]{Lemma}
\newtheorem{corollary}[theorem]{Corollary}
\newtheorem{definition}{Definition}

\lstset{
  language=C++,
  basicstyle=\ttfamily\small,
  keywordstyle=\bfseries,
  breaklines=true,
  numbers=left,
  numberstyle=\tiny,
  frame=single,
  tabsize=2
}

\title{IOI 2023 --- Soccer Stadium}
\author{}
\date{}

\begin{document}
\maketitle

\section{Problem Statement}

Given an $N \times N$ grid where some cells contain trees ($F[r][c] = 1$),
find the largest \emph{regular stadium}: a set of empty cells such that
any two cells can be connected by at most two straight kicks
(horizontal or vertical line segments, with every intermediate cell in
the stadium).

\section{Characterization}

\begin{definition}
A \emph{regular stadium} is a set $\mathcal{S}$ of empty cells such that
for every pair of cells $(r_1, c_1), (r_2, c_2) \in \mathcal{S}$, there
exists a cell $(r_m, c_m) \in \mathcal{S}$ with:
\begin{itemize}
  \item the entire row segment from $(r_1, c_1)$ to $(r_1, c_m)$ (or the
        column segment from $(r_1, c_1)$ to $(r_m, c_1)$) lies in
        $\mathcal{S}$, and
  \item the entire column (or row) segment from $(r_m, c_m)$ to
        $(r_2, c_2)$ lies in $\mathcal{S}$.
\end{itemize}
\end{definition}

\begin{theorem}[Column-interval relay]\label{thm:relay}
Every optimal regular stadium can be described as follows: there exists
a contiguous column interval $[c_l, c_r]$ (the \emph{relay zone}) such
that the stadium consists of cells from rows that contain no tree in
$[c_l, c_r]$, extended left and right as far as possible in each row.

Formally, for each row $r$ where $F[r][c] = 0$ for all $c \in [c_l, c_r]$,
the stadium includes the maximal contiguous empty segment in row $r$
that contains $[c_l, c_r]$.
\end{theorem}

\begin{proof}[Proof sketch]
In a regular stadium, any two cells must connect via at most two kicks.
Consider the set of columns that every included row-segment covers.
The intersection of all row-segments' column ranges must be non-empty
(otherwise two cells in rows with disjoint column ranges cannot connect
in two kicks).  This common column interval serves as the relay zone.

Any two cells $(r_1, c_1)$ and $(r_2, c_2)$ can then connect via:
kick~1 along row $r_1$ to some column $c \in [c_l, c_r]$, then kick~2
along column $c$ to row $r_2$, then (if needed) kick along row $r_2$ to
$c_2$ --- but since both rows span $[c_l, c_r]$, the column kick from
$(r_1, c)$ to $(r_2, c)$ must pass through only rows included in the
stadium.

For the column segment to be fully inside the stadium, every row between
$r_1$ and $r_2$ must include column $c$.  This means we can only include
a contiguous vertical block of rows (those with $[c_l, c_r]$ clear).
More precisely, the included rows need not be contiguous: we need every
row between any two included rows (at column $c$) to also be included.
So the set of included rows, restricted to any relay column, must be
contiguous.

Actually, we only need \emph{some} relay column $c \in [c_l, c_r]$ such
that the column of included rows at $c$ covers all row-pairs.  Since the
included rows are exactly those where $[c_l, c_r]$ is tree-free, and
between any two such rows a row with a tree in $[c_l, c_r]$ would block
the column kick, we must ensure the included rows form contiguous blocks
per relay column.  The simplest valid structure is: enumerate all
$(c_l, c_r)$ pairs, for each find the maximal set of contiguous rows
where $[c_l, c_r]$ is tree-free, and extend each row as far as possible.
\end{proof}

\begin{corollary}
The optimal stadium size can be found by enumerating all $O(N^2)$
column-interval pairs $(c_l, c_r)$ and computing the total cells for
each in $O(N)$ time.
\end{corollary}

\section{Algorithm}

\begin{enumerate}
  \item Precompute, for each cell $(r, c)$, the leftmost and rightmost
        extent of the contiguous empty segment in row $r$ containing $c$.
  \item Precompute row prefix sums to check whether $[c_l, c_r]$ is
        tree-free in row $r$ in $O(1)$.
  \item For each pair $(c_l, c_r)$ with $0 \le c_l \le c_r < N$:
        \begin{enumerate}
          \item For each row $r$: if $[c_l, c_r]$ is tree-free, find the
                maximal empty segment containing $[c_l, c_r]$ and add its
                length.
          \item The rows must form contiguous blocks for the column kicks
                to work.  We sum over all rows in each maximal contiguous
                block and take the maximum block-sum, or --- since the
                relay allows reaching any row through the relay columns
                --- we need all included rows to be mutually reachable.
                The simplest correct approach: only include rows that form
                one contiguous block (no tree row in between at any relay
                column), or sum all qualifying rows and verify connectivity.
        \end{enumerate}
  \item Additionally check single-row and single-column stadiums, and
        the maximal-rectangle baseline.
\end{enumerate}

In fact, the simplest correct $O(N^3)$ approach that handles all cases:
for each $(c_l, c_r)$, sum the extended row lengths over \emph{all}
qualifying rows.  Two cells in different qualifying rows can always
connect: go along row to some column $c \in [c_l, c_r]$, then down
column $c$ to the target row (all intermediate rows also qualify since
they must be tree-free in $[c_l, c_r]$ for this to work).  So we also
need the qualifying rows to be contiguous.  We handle this by finding
maximal contiguous runs.

\section{Implementation}

\begin{lstlisting}
#include <bits/stdc++.h>
using namespace std;

int biggest_stadium(int N, vector<vector<int>> F) {
    // Precompute: for each (r, c), the left/right extent of empty segment
    vector<vector<int>> lft(N, vector<int>(N)), rgt(N, vector<int>(N));
    for (int r = 0; r < N; r++) {
        for (int c = 0; c < N; c++)
            lft[r][c] = (F[r][c] == 1) ? c : (c > 0 ? lft[r][c-1] : 0);
        for (int c = N - 1; c >= 0; c--)
            rgt[r][c] = (F[r][c] == 1) ? c : (c < N-1 ? rgt[r][c+1] : N-1);
    }

    // Row prefix sums of trees
    vector<vector<int>> rpfx(N, vector<int>(N + 1, 0));
    for (int r = 0; r < N; r++)
        for (int c = 0; c < N; c++)
            rpfx[r][c+1] = rpfx[r][c] + F[r][c];

    auto row_clear = [&](int r, int cl, int cr) -> bool {
        return rpfx[r][cr+1] - rpfx[r][cl] == 0;
    };

    int best = 0;

    // Enumerate relay column intervals [cl, cr]
    for (int cl = 0; cl < N; cl++) {
        for (int cr = cl; cr < N; cr++) {
            // Find contiguous runs of rows where [cl, cr] is tree-free.
            // For each run, sum the extended row lengths.
            int run_sum = 0;
            for (int r = 0; r < N; r++) {
                if (row_clear(r, cl, cr)) {
                    // Extend row r: leftmost empty from cl, rightmost from cr
                    int left = lft[r][cl];
                    int right = rgt[r][cr];
                    run_sum += right - left + 1;
                } else {
                    // Tree in [cl, cr] at row r: end of contiguous run
                    best = max(best, run_sum);
                    run_sum = 0;
                }
            }
            best = max(best, run_sum);
        }
    }

    return best;
}
\end{lstlisting}

\section{Complexity}

\begin{itemize}
  \item \textbf{Time:} $O(N^3)$.  There are $O(N^2)$ column-interval
        pairs; for each, a single pass over $N$ rows.
  \item \textbf{Space:} $O(N^2)$ for the precomputed arrays.
\end{itemize}

\begin{theorem}[Correctness]
The algorithm correctly finds the largest regular stadium.
\end{theorem}

\begin{proof}
By Theorem~\ref{thm:relay}, every regular stadium has a relay column
interval $[c_l, c_r]$ such that each row's contribution is a contiguous
empty segment spanning $[c_l, c_r]$, and the contributing rows form a
contiguous vertical block (so that column kicks between any two rows
stay within the stadium).

Our algorithm enumerates all possible relay intervals and, for each,
finds the best contiguous vertical block of qualifying rows.  Within
such a block, any two cells connect in at most two kicks: a horizontal
kick to a relay column, then a vertical kick to the target row (and
optionally a horizontal kick in the target row).

Since we maximize over all choices, the algorithm finds the global
optimum.
\end{proof}

\end{document}
